\section{Future work}\label{sec:future-work}


Theorem~\ref{thm:scoped-exact-six} and the limitations of
Section~\ref{sec:limitations-nonclaims} suggest several follow-on
research directions. Each subsection below states a direction, not a
commitment.

\subsection{Endogenous-instrument meta-limit}\label{subsec:future-meta-limit}

A formalization of the meta-limit described in
Section~\ref{subsec:meta-limit} would first define internal arguments and the soundness and completeness of
an instrument family, and then target a precise form of the heuristic of
Section~\ref{subsec:meta-limit}: that no internal argument can simultaneously
establish the scoped exact-six claim, certify the soundness of the
accepted instrument family, and certify its completeness for
exact-six judgments formulable inside the same regime.
Theorem~\ref{thm:scoped-exact-six} would remain intact under such a
result; the additional content would be a precise account of
self-certification boundaries. A closely related open question
concerns recognition. Not every finite audited feature of a description in
\FATCD{} is recognized: the bare normalization predicate of
Remark~\ref{rem:recognition-open} receives the unresolved label (x). What remains open is
which classes of assertion-bearing records --- normalization constraints not
presented as the law of a transition record (for instance for signed weights),
finite conservation laws, emergent composite obstructions --- admit a reduction
in which the assertion is itself a required field of an active cluster
(label (i)) without adjoining wrapper records such as a realization relation
and a one-class packaging map (Remark~\ref{rem:recognition-open}), or a trace
reference, and whether any of
them carries the recorded witness of category (xi) of
Definition~\ref{def:residual-scheme}, which is a recorded screen rather than
the identification of a new closure role.

\subsection{Broader-domain exact-six}\label{subsec:future-broader-domain}

Candidate extensions beyond \FATCD{} include admissibility regimes
that weaken finiteness, closure, or audit
(Section~\ref{sec:fatcd-domain}), and regimes that admit annotation
kinds beyond the six \(\mathrm{Ann}_*\) categories of
Definition~\ref{def:annotations}. Any such extension would require
fresh upper-bound and decomposition/exhaustion analysis, since the
alignment rules and threat closures are not stated at scope-free
strength.

\subsection{Uniqueness or canonical decomposition}\label{subsec:future-uniqueness}

Theorem~\ref{thm:scoped-exact-six} is existential. A uniqueness study
would investigate when distinct well-formed decomposition/exhaustion
records of the same \(D \in \FATCD{}\) are equivalent under explicit
translation records, or when a canonical representative can be
selected. The channel-versus-activation discipline
(Proposition~\ref{prop:channel-vs-activation}) is compatible with
such a study; no uniqueness claim is part of the present theorem.

\subsection{Full six-primitives algebra}\label{subsec:future-algebra}

The generator-and-relation structure of \Pone{}--\Psix{} beyond
typed non-collapse remains open. A full algebra would require
scope-free independence data, not just the local boundary
distinctions of Table~\ref{tab:boundary-matrix}, and would likely
require a categorical apparatus richer than the raw grammar of
Section~\ref{subsec:raw-grammar}.

\subsection{Empirical substrate mapping}\label{subsec:future-substrate}

The empirical-bridge admissibility regime
(Definition~\ref{prop:empirical-bridge}) can be connected to
concrete substrate-level observables. The present paper treats
empirical bridges as threshold-dependent, level-indexed, and
instrument-visible annotations; later work can couple this to
substrate-specific measurement theory while preserving the nonclaim
of empirical realization.

\subsection{Mechanization}\label{subsec:future-mechanization}

A Lean~4 development accompanies the present paper;
Appendix~\ref{app:mechanization} records, claim by claim, how far each definition
and theorem is currently mechanized. Strengthening that mechanization toward full
proofs --- in particular a Lean account of the boundary distinctions, the
decomposition-existence construction, and the recognition hypothesis underlying
the upper-bound classification --- is a natural direction, building on the
closure-related infrastructure already mechanized in the Foundations
program~\citep{Tsiokos2026SixBirdsFoundations,DeMouraUllrich2021}.

\subsection{Cross-instrument comparison}\label{subsec:future-cross-instrument}

Decomposition/exhaustion records of the same \(D \in \FATCD{}\) can
be compared under different admissible instrument assignments.
Definition~\ref{prop:visibility} admits visibility records that
suppress different portions of a claim under different instruments; a
cross-instrument study could make the instrument-dependence of
\(\Dec{D}\) explicit and sharpen the notion of instrument-relative
claim strength.
